\section{Spectral properties of bounded self-adjoint linear operators}

1
2 Callback to the finite spectral theorem! 
4
5 Recall from Axler that S = W * TW is equivalence of matrices. They are intuitively the same thing "up to a change of orthonormal basis". So naturally we wish self-adjointnes to be preserved
6
7 Extra: give an exact criterion for when such a diagonal operator is self-adjoint.
8
9 Notice that operators like this would be the analogue of diagonal matrices. So we're definitely not in Kansas anymore, not even "diagonal operators" would have eigenvalues! Gone is the nice spectral theorem we got in the finite-dim or compact case.
10 This is a very famous result called the Hellinger-Toeplitz theorem. It is not a nice result, it actually shows that quantum mechanics is inherently very difficult. Self-adjointness is something we really do want in QM, but most QM operators tend not to be bounded (bounded operators actually are the experiments whose outcomes are a bounded set). So the only thing we can do is to take operators which are not everywhere defined! Because of this theorem, the entire chapter of Kreyszig on unbounded operators exists and is immensely important in physics! I recommend you to check out the wiki of Helling-Toeplitz for more information

\begin{exercise}{1}
It was mentioned in the text that for all self-adjoint linear operator $T$ the inner product $\brackets{Tx, x}$ is real.
What does this imply for matrices?
What familiar theorem on matrices does Theorem 9.1-1 include as a special case?
\end{exercise}
\begin{proof}
fill
\end{proof} 

\begin{exercise}{2}
If in the finite dimensional case, a self-adjoint linear operator $T$ is represented by a diagonal matrix, show that the matrix must be real.
What is the spectrum of $T$?
\end{exercise}
\begin{proof}
fill
\end{proof} 

\begin{exercise}{4}
Illustrate Theorem 9.1-2 by an operator whose spectrum consists of a single given value $\lambda_0$.
What is the largest $c$ in this case?
\end{exercise}
\begin{proof}
fill
\end{proof} 

\begin{exercise}{5}
Let $T: H \to H$ and $W: H \to H$ be bounded linear operators on a complex Hilbert space $H$.
If $T$ is self-adjoint, show that $S = W^\ast T W$ is self-adjoint.
\end{exercise}
\begin{proof}
fill
\end{proof} 

\begin{exercise}{6}
Let $T: \ell^2 \to \ell^2$ be defined by $(x_1, x_2, \dots) \mapsto (0, 0, x_1, x_2, \dots)$.
Is $T$ bounded?
Self-adjoint?
Find $S: \ell^2 \to \ell^2$ such that $T = S^2$.
\end{exercise}
\begin{proof}
fill
\end{proof} 

\begin{exercise}{7}
Let $T: \ell^2 \to \ell^2$ be defined by $y = (y_n) = T x$, $x = (x_n)$, $y_n = \lambda_n x_n$, where $(\lambda_n)$ is a bounded sequence on $\R$, and $a = \inf \lambda_j, b = \sup \lambda_j$.
Show that each $\lambda_j$ is an eigenvalue of $T$.
Under what condition will $[a, b] \subseteq \sigma(T)$.
\end{exercise}
\begin{proof}
fill
\end{proof} 

\begin{exercise}{8}
Using Theorem 9.1-2, prove that the spectrum of the operator $T$ in exercise 7 is the closure of the set of eigenvalues.
\end{exercise}
\begin{proof}
fill
\end{proof} 

\begin{exercise}{9}
From 2.2-7 and 3.1-5 we know that the Hilbert space $L^2 [0, 1]$ is the completion of the inner product space $X$ of all continuous functions on $[0, 1]$ with inner product defined by $\brackets{x, y} = \int_0^1 x(t) \overline{y(t)} dt$.
Show that $T: L^2 [0, 1] \to L^2 [0, 1]$ defined by $y(t) = T x(t) = t x(t)$ is a bounded self-adjoint linear operator without eigenvalues.
\end{exercise}
\begin{proof}
fill
\end{proof} 

\begin{exercise}{10}
It is interesting to note that a linear operator $T$ which is defined on all of a complex Hilbert space $H$ and satisfies $\brackets{Tx, y} = \brackets{x, Ty}$ for all $x, y \in H$ must be bounded (so that our assumption of boundedness at the beginning of the text would not be necessary).
Prove this fact.
\end{exercise}
\begin{proof}
fill
\end{proof} 
